% GENERATED FILE. Edit canonical lessons, then run npm run generate:books.
% canonical-source-hash: fnv1a64:ffd738d1611cc2bf
% canonical-lessons: MR-C95-shabda, MR-C95-khel, MR-C95-patra, MR-C95-chitra, MR-C95-jhad

\chapter{Word, Game, Letter, Picture, Tree}
\label{ch:mr-word-game-letter-picture-tree}
\hlchaptermodality{95}

\begin{chapteropening}
\textbf{By the end of this chapter:} \emph{I can say a word, a game, a letter, a picture and a tree.}

Complete the last lesson of chapter 95: i can say a word, a game, a letter, a picture and a tree.
\end{chapteropening}

\section[shabda]{\mr{शब्द} (shabda) — a word}

\label{lesson:MR-C95-shabda}



\begin{quote}
Before the new one: say the Marathi for a question, then the Marathi for an answer.
\end{quote}

\subsection*{You'll want to know: \mr{शब्द}}
\textbf{\mr{शब्द}} — \emph{shabda} — \textquotedblleft{}a word\textquotedblright{}.

\begin{grammarlens}[title={What you've built}]
One more word for this chapter.
\end{grammarlens}

\subsection*{Your turn}
\begin{itemize}
\raggedright
  \item \emph{Say it:} \emph{shabda}
  \item \emph{Say it:} \emph{shabda}, once more
  \item \emph{Say it:} say \emph{shabda}
  \item \emph{Recall:} say the Marathi for a question, then the Marathi for an answer, then say \emph{shabda} again
\end{itemize}

\begin{tcolorbox}[breakable,colback=teal!4,colframe=teal!35!black,title={Before you move on}]
What does \mr{शब्द} mean? (\textquotedblleft{}A word\textquotedblright{}.) Say it once more.
\end{tcolorbox}
\section[kheḷ]{\mr{खेळ} (kheḷ) — a game}

\label{lesson:MR-C95-khel}



\begin{quote}
Before the new one: say the Marathi for an answer, then the Marathi for a word.
\end{quote}

\subsection*{You'll want to know: \mr{खेळ}}
\textbf{\mr{खेळ}} — \emph{kheḷ} — \textquotedblleft{}a game\textquotedblright{}.

\begin{grammarlens}[title={What you've built}]
One more word for this chapter.
\end{grammarlens}

\subsection*{Your turn}
\begin{itemize}
\raggedright
  \item \emph{Say it:} \emph{kheḷ}
  \item \emph{Say it:} \emph{kheḷ}, once more
  \item \emph{Say it:} say \emph{kheḷ}
  \item \emph{Recall:} say the Marathi for an answer, then the Marathi for a word, then say \emph{kheḷ} again
\end{itemize}

\begin{tcolorbox}[breakable,colback=teal!4,colframe=teal!35!black,title={Before you move on}]
What does \mr{खेळ} mean? (\textquotedblleft{}A game\textquotedblright{}.) Say it once more.
\end{tcolorbox}
\section[patra]{\mr{पत्र} (patra) — a letter (post)}

\label{lesson:MR-C95-patra}



\begin{quote}
Before the new one: say the Marathi for a word, then the Marathi for a game.
\end{quote}

\subsection*{You'll want to know: \mr{पत्र}}
\textbf{\mr{पत्र}} — \emph{patra} — \textquotedblleft{}a letter (post)\textquotedblright{}.

\begin{grammarlens}[title={What you've built}]
One more word for this chapter.
\end{grammarlens}

\subsection*{Your turn}
\begin{itemize}
\raggedright
  \item \emph{Say it:} \emph{patra}
  \item \emph{Say it:} \emph{patra}, once more
  \item \emph{Say it:} say \emph{patra}
  \item \emph{Recall:} say the Marathi for a word, then the Marathi for a game, then say \emph{patra} again
\end{itemize}

\begin{tcolorbox}[breakable,colback=teal!4,colframe=teal!35!black,title={Before you move on}]
What does \mr{पत्र} mean? (\textquotedblleft{}A letter (post)\textquotedblright{}.) Say it once more.
\end{tcolorbox}
\section[chitra]{\mr{चित्र} (chitra) — a picture}

\label{lesson:MR-C95-chitra}



\begin{quote}
Before the new one: say the Marathi for a game, then the Marathi for a letter.
\end{quote}

\subsection*{You'll want to know: \mr{चित्र}}
\textbf{\mr{चित्र}} — \emph{chitra} — \textquotedblleft{}a picture\textquotedblright{}.

\begin{grammarlens}[title={What you've built}]
One more word for this chapter.
\end{grammarlens}

\subsection*{Your turn}
\begin{itemize}
\raggedright
  \item \emph{Say it:} \emph{chitra}
  \item \emph{Say it:} \emph{chitra}, once more
  \item \emph{Say it:} say \emph{chitra}
  \item \emph{Recall:} say the Marathi for a game, then the Marathi for a letter, then say \emph{chitra} again
\end{itemize}

\begin{tcolorbox}[breakable,colback=teal!4,colframe=teal!35!black,title={Before you move on}]
What does \mr{चित्र} mean? (\textquotedblleft{}A picture\textquotedblright{}.) Say it once more.
\end{tcolorbox}
\section[jhāḍ]{\mr{झाड} (jhāḍ) — a tree}

\label{lesson:MR-C95-jhad}



\begin{quote}
Before the new one: say the Marathi for a letter, then the Marathi for a picture.
\end{quote}

\subsection*{You'll want to know: \mr{झाड}}
\textbf{\mr{झाड}} — \emph{jhāḍ} — \textquotedblleft{}a tree\textquotedblright{}.

\begin{grammarlens}[title={What you've built}]
One more word for this chapter.
\end{grammarlens}

\subsection*{Your turn}
\begin{itemize}
\raggedright
  \item \emph{Say it:} \emph{jhāḍ}
  \item \emph{Say it:} \emph{jhāḍ}, once more
  \item \emph{Say it:} say \emph{jhāḍ}
  \item \emph{Recall:} say the Marathi for a letter, then the Marathi for a picture, then say \emph{jhāḍ} again
\end{itemize}

\begin{tcolorbox}[breakable,colback=teal!4,colframe=teal!35!black,title={Before you move on}]
What does \mr{झाड} mean? (\textquotedblleft{}A tree\textquotedblright{}.) Say it once more.
\end{tcolorbox}
